% TOP_PROBLEM: {"id":30007561,"problem_number":"LOCAL-30007561","title":"Welfare with limited preference information","category_id":21,"category":{"id":21,"name":"economics","display_name":"Economics","description":"Open economic theory statements, published research agendas, and proposed scoped specifications; question provenance and historical priority are qualified per record.","slug":"economics","order_index":21,"created_at":"2026-10-08T00:00:00Z"},"status":"research_question","external_url":null,"legacy_ids":[],"published":false,"statement":"Determine optimal welfare distortion and query policies when normalized utilities are initially known only ordinally and the mechanism can elicit a limited number of exact cardinal values.","economics":{"original_id":50,"field":"Economic theory, equilibrium and social choice","kind":"theoretical","formalization_basis":"adapted_research_question","source_status":"research_agenda","priority_status":"earliest_found","priority_note":"This is the earliest source in which the relevant agenda was verified during this audit; absolute priority has not been established.","earliest_verified_reference":{"authors":["Elliot Anshelevich","Aris Filos-Ratsikas","Nisarg Shah","Alexandros A. Voudouris"],"title":"Distortion in Social Choice Problems: The First 15 Years and Beyond","year":2021,"url":"https://arxiv.org/pdf/2103.00911","doi":"10.24963/ijcai.2021/589","locator":"§5 Generalizations of Distortion, Limited cardinal queries, Open Problems 3–4, PDF p.9","evidence_summary":"The stated agenda asks which settings admit query–distortion tradeoffs and how limited cardinal information can improve distortion."},"current_reference":{"authors":["Elliot Anshelevich","Aris Filos-Ratsikas","Nisarg Shah","Alexandros A. Voudouris"],"title":"Distortion in Social Choice Problems: The First 15 Years and Beyond","year":2021,"url":"https://arxiv.org/pdf/2103.00911","doi":"10.24963/ijcai.2021/589","locator":"§5 Generalizations of Distortion, Limited cardinal queries, Open Problems 3–4, PDF p.9","evidence_summary":"The stated agenda asks which settings admit query–distortion tradeoffs and how limited cardinal information can improve distortion."},"formalization_note":"This is a scoped family adapted from Open Problems3–4. Individual parameter regimes need a current literature check before promotion to explicit_open."},"statusQualification":"Published question or research agenda verified at the cited locator. The mathematical specification is adapted for this catalogue and includes additional assumptions; source publication does not establish first-ever priority or certify this exact model remains unsolved. This is a scoped family adapted from Open Problems3–4. Individual parameter regimes need a current literature check before promotion to explicit_open.","statusReviewedAt":"2026-10-08"}
% FIRST_POSTED: 2026-10-08
% ENTRY_KIND: statement_only
% !TeX program = lualatex
\documentclass[11pt]{article}
\usepackage{fontspec}
\usepackage[a4paper,margin=25mm]{geometry}
\usepackage{amsmath,amssymb,mathtools}
\usepackage{xurl}
\usepackage[hidelinks]{hyperref}
\newcommand{\sref}[2]{\hyperlink{source:#1:#2}{[S#2]}}
\setlength{\emergencystretch}{3em}
\begin{document}
% problemId:30007561
\section*{Welfare with limited preference information}\label{30007561}
\subsection{Definitions and mathematical statement}
Let \(n\) voters choose one of \(m\geq2\) alternatives. Voter \(i\) has unknown nonnegative utilities \(u_i(a)\) normalized by \(\sum_{a=1}^m u_i(a)=1\). The mechanism initially observes each voter's ordinal ranking, with a fixed rule for utility ties, and may adaptively query at most \(k\) utility values per voter. Queries return truthful exact values; strategic reporting is outside this specification. A possibly randomized algorithm \(A\) selects an alternative, and social welfare is \(W_u(a)=\sum_i u_i(a)\). Define its worst-case distortion by
\[
 D(A)=\sup_u\frac{\max_a W_u(a)}{E[W_u(A(u))]},\qquad
 D^*(n,m,k)=\inf_A D(A),
\]
where the supremum ranges over normalized utility profiles and the expectation over algorithmic randomization. Determine sharp orders and constants for \(D^*(n,m,k)\) in parameter regimes not covered by existing results, together with optimal adaptive query policies and matching information-theoretic lower bounds. Explicitly state the regime and all known upper and lower bounds before claiming a residual gap. Examine whether bounded measurement error or a restricted nonnegative valuation domain changes the query–welfare frontier. The verified source supplies broader numbered open questions about cardinal-information tradeoffs; it does not certify every special case of this proposed frontier as unresolved, and solved ordinal-distortion conjectures are not reinstated.

\par\textbf{Formulation and scope.} This is a scoped family adapted from Open Problems3–4. Individual parameter regimes need a current literature check before promotion to explicit\_open.

\par\textbf{Open-question provenance.} An explicit question or research agenda was verified at the source locator. This does not by itself certify that every later version remains unresolved \sref{30007561}{1}.

\par\textbf{Historical priority.} This is the earliest source in which the relevant agenda was verified during this audit; absolute priority has not been established.

\subsection{Short English statement}
Determine optimal welfare distortion and query policies when normalized utilities are initially known only ordinally and the mechanism can elicit a limited number of exact cardinal values.

\subsection{Sources}
\begin{itemize}\raggedright
\item[\textnormal{[S1]}]\hypertarget{source:30007561:1}{} Elliot Anshelevich, Aris Filos-Ratsikas, Nisarg Shah, Alexandros A. Voudouris. 2021. Distortion in Social Choice Problems: The First 15 Years and Beyond. §5 Generalizations of Distortion, Limited cardinal queries, Open Problems 3–4, PDF p.9.
\url{https://arxiv.org/pdf/2103.00911}.
DOI: 10.24963/ijcai.2021/589.
{\small\textit{Earliest verified question/agenda reference; historical priority qualified below. The stated agenda asks which settings admit query–distortion tradeoffs and how limited cardinal information can improve distortion.}}
\end{itemize}
\end{document}
